\section{Metriche di rete}
\subsection{Architettura di rete a strati}
\subsubsection*{Concetti di livello, servizio, interfaccia e protocollo di un'architettura a strati}
Un'architettura a strati è un modello concettuale che suddivide il funzionamento di un sistema in una pila gerarchica
di livelli o strati. Ogni livello svolge una parte di compiti, necessari il funzionamento corretto dell'intero sistema.
Permette di suddividere un problema complesso in più sotto-problemi più semplici da risolvere, favorendo la modularità,
la riusabilità e la manutenibilità del sistema.

Nelle implementazioni reali si cerca di mantenere il più possibile indipendenti i livelli, in modo da garantire modularità
e riusabilità del sistema, ma in alcuni casi si utilizzano implementazioni cross-layer per migliorare l'efficienza e le
prestazioni del sistema.

Un'architettura a strati si basa su quattro concetti fondamentali:
\begin{itemize}
	\item \textbf{livello}: insieme di componenti che risolvono una serie di sottoproblemi omogenei del problema complesso;
	ogni livello fornisce un insieme di servizi al livello superiore nascondendone i dettagli di implementazione
	\item \textbf{servizio}: insieme di primitive o funzionalità che un livello offre al livello superiore
	\item \textbf{interfaccia}: insieme di regole e convenzioni che definiscono come un livello può accedere ai servizi
	del livello sottostante
	\item \textbf{protocollo}: insieme di regole e convenzioni che definiscono come due livelli corrispondenti situati
	in nodi diversi della rete comunicano tra loro
\end{itemize}

\subsubsection*{Livelli nelle reti di calcolatori - modello ISO/OSI}
Il modello ISO/OSI (Open Systems Interconnection) è il modello di riferimento per le reti di calcolatori,
suddivide il funzionamento di una rete nei seguenti sette livelli:
\begin{itemize}
	\item[0.] \textbf{physical medium}: rappresenta il mezzo fisico di trasmissione di segnali analogici, come cavi, fibre
	ottiche, onde radio, ecc., collega fisicamente i vari nodi della rete
	\item[L1.] \textbf{physical layer}: fornisce un canale digitale di trasmissione a bit inaffidabile tra due nodi adiacenti della
	rete, gestisce la conversione digitale/analogica e viceversa per trasmettere i bit digitali sul mezzo fisico tramite
	segnali analogici
	\item[L2.] \textbf{data link}: fornisce un canale digitale di trasmissione a bit affidabile tra due nodi adiacenti della rete,
	gestisce la correzione di errori sul bit introdotti dal canale inaffidabile del livello data link
	\item[L3.] \textbf{network}: fornisce un canale digitale di trasmissione a pacchetti inaffidabile tra due nodi non
	necessariamente adiacenti della rete, gestisce l'instradamento dei pacchetti attraverso la rete (non garantisce che i
	pacchetti arrivino a destinazione, né che arrivino in ordine)
	\item[L4.] \textbf{transport}: fornisce un canale di comunicazione end-to-end tra due processi in esecuzione su due nodi
	doversi della rete, gestisce la segmentazione dei dati in pacchetti, la correzione di errori sul pacchetto e il controllo
	del flusso di trasmissione dei pacchetti (ad esempio l'avvenuta ricezione)
	\item[L5.] \textbf{session}: fornisce una connessione logica tra due processi in esecuzione su due nodi diversi della
	rete, gestisce la creazione e la distruzione di un canale di comunicazione tra due processi
	\item[L6.] \textbf{presentation}: fornisce un sistema di scambio dei dati tra due processi in esecuzione su due nodi
	diversi della rete, gestisce la codifica, compressione, cifratura e formattazione dei dati da trasmettere
	\item[L7.] \textbf{application}: fornisce un'interfaccia tra l'utente e la rete, gestisce gli input dell'utente e la
	 visualizzazione dei dati ricevuti dalla rete
\end{itemize}

\subsubsection*{Incapsulamento dei dati e aggiunta degli headers}
In fase di trasmissione, ogni livello di rete, riceve i dati dal livello superiore detti payload e, prima di trasmetterli
al livello inferiore, ci aggiunge un header contenente informazioni di protocollo necessarie al livello corrispondente del
nodo ricevente per poter interpretare correttamente i dati ricevuti. Il pacchetto formato da header e payload ad un certo
livello diventa, quindi, il payload del livello inferiore.

Viceversa in fase di ricezione, ogni livello di rete, riceve i dati dal livello inferiore, separa l'header di livello
dal payload, processa il payload con le informazioni dell'header appena estratto e passa i dati processati al livello
superiore senza l'header di livello appena consumato.

\newpage

\subsection{Paradigma di trasmissione a pacchetto}
\subsubsection*{Pacchettizzazione del livello 4 di transport}
Il livello 4 di transport suddivide i dati da trasmettere, ricevuti dai livello superiori, in pacchetti per poterli
trasmettere agevolmente attraverso la rete. Teoricamente la pacchettizzazione dovrebbe essere invisibile ai livelli
superiori però in alcuni casi è utile attuare un'implementazione cross-layer in cui i pacchetti coincidono le code
unit dei codec video (blocchi in cui vengono divisi i frame).

\subsubsection*{Vantaggi della trasmissione a pacchetti}
I vantaggi della trasmissione a pacchetti sono:
\begin{itemize}
	\item \textbf{efficienza}: i pacchetti possono essere instradati in modo indipendente attraverso la rete,
	permettendo di sfruttare al meglio le risorse disponibili e ridurre i tempi di trasmissione
	\item \textbf{affidabilità}: in caso di errore o perdita di pacchetti, è possibile ritrasmettere solo i
	pacchetti persi o danneggiati, invece di dover ritrasmettere l'intero flusso di dati
	\item \textbf{scalabilità}: la trasmissione a pacchetti consente di gestire la comunicazione tra un numero
	elevato di nodi e dispositivi con grandi volumi di traffico senza cambiare l'architettura dell'infrastruttura
	\item \textbf{multiplexing statistico}: più flussi di dati possono condividere le stesse risorse di rete
	garantendo una maggiore flessibilità
\end{itemize}

\subsection{Metriche di traffico}
\subsubsection*{Unità di misura per il tasso di trasmissione}
Il tasso di trasmissione o bitrate si può misura in:
\begin{itemize}
	\item bit/s (bps): bit per secondo, con relativi multipli kbps, Mbps, Gbps, Tbps
	\item byte/s (Bps): byte per secondo, con relativi multipli kBps, MBps, GBps, TBps, con 1 B = 8 b
\end{itemize}
Per convenzione, non si utilizza la notazione binaria per cui 1 MB/s =\(2^{20}\) B/s, bensì si usa quella decimale in cui
1 MB/s = \(10^6\) B/s.

\subsubsection*{Bitrate, throughput, goodput}
Tutti e tre i termini indicano il tasso o velocità di trasmissione dei dati tra due entità peer (livelli corrispondenti)
di due nodi della rete, ma differiscono per il livello in cui sono misurati:
\begin{itemize}
	\item \textbf{bitrate}: velocità nominale del livello fisico L1, indicato con \(R_0\), dipende dalla larghezza di
	banda, dalla modulazione utilizzata e dal rapporto segnale/rumore del canale, è il tasso massimo di trasmissione
	dei bit su canale single-hop inaffidabile (canale tra nodi adiacenti senza garanzia di correttezza e consegna dei bit)
	\item \textbf{throughput}: velocità media misurata tra il livello 2 di data link e il livello 4 di transport,
	indicata con \(S\), in base al livello in cui si misura può tenere conto dell'overhead dovuto agli error
	correction codes, agli header o alla ritrasmissione dei pacchetti persi o danneggiati
	\item \textbf{goodput}: velocità effettiva misurata in un livello superiore al 4 (session, presentation o application),
	corrisponde al throughput a lungo termine \(S_n\) e rappresenta il tasso di trasmissione dei dati utili, tenendo conto di tutti
	gli overhead necessari alla trasmissione dei dati
\end{itemize}
In generale si ha che \(R_0 \geq S \geq \text{goodput}\), ovvero il tasso di trasmissione diminuisce con l'aumentare dei
livelli di rete a causa dell'overhead introdotto dagli header di ogni livello.

\subsubsection*{Protocol Data Unit, Protocol Control Information e Service Data Unit}
\begin{itemize}
	\item \textbf{Protocol Data Unit (PDU)}: insieme di dati che viene trasmesso tra due livelli corrispondenti di due nodi
	della rete, coincide con il payload, ovvero i dati utili da trasmettere
	\item \textbf{Protocol Control Information (PCI)}: insieme di informazioni di protocollo contenute nell'header di un
	livello, costituisce l'overhead introdotto dal livello
	\item \textbf{Service Data Unit (SDU)}: insieme di dati che viene passato al livello inferiore, è composta dal payload
	e dall'header del livello superiore, e diventa il payload o SDU del livello inferiore.
\end{itemize}

\subsubsection*{Definizioni matematiche di throughput e goodput}
Si considera \(b_n(t)\) come il numero di bit trasmessi nell'intervallo di tempo \([0,t]\) dal livello \(n\) di rete
al livello inferiore, e si definiscono le seguenti metriche di throughput e goodput:
\[\begin{array}{c c}
	\displaystyle S_n(t;T) = \frac{b_n(t) - b_n(t-T)}{T} \quad\, & \quad \text{throughput medio nell'intervallo di tempo } [t,t+T] \\[10pt]
	\displaystyle S_n(t) = \frac{d \, b_n(t)}{d \, t} \qquad\qquad\;\, ~ & \quad \text{throughput istantaneo al tempo } t \\[12pt]
	\displaystyle S_n = \lim_{t,T \to \infty} S_n(t;T)  & \quad \text{goodput o throughput a lungo termine}
\end{array}\]

\noindent
Si definiscono inoltre le metriche di throughput e l'efficienza in funzione della taglia delle PDU e SDU:
\[\begin{array}{c c}
	\displaystyle S_{n+1} = S_n \cdot \frac{|SDU_n|}{|PDU_n|} = S_n \cdot \frac{|SDU_n|}{|PCI_n| + |SDU_n|} & \begin{array}{c} \text{throughput massimo del} \\ \text{livello } n+1 \text{ in funzione} \\ \text{di quello del livello } n \end{array} \\[20pt]
	\displaystyle \eta_n = \frac{S_{n+1}}{S_n} = \frac{|SDU_n|}{|PDU_n|} = \frac{|SDU_n|}{|PCI_n| + |SDU_n|} = \frac{1}{1 + \frac{|PCI_n|}{|SDU_n|}} = \frac{|PDU_{n+1}|}{|PDU_n|} & \begin{array}{c} \text{efficienza del livello } n \\ \text{con } 0 \leq \eta_n \leq 1 \end{array}
\end{array}\]

\noindent
L'efficienza di un livello \(n\) indica la porzione di throughput del livello \(n\) che viene effettivamente utilizzata
per trasmettere i dati utili del livello \(n+1\).

\subsubsection*{Bitrate di una connessione multi-hop}
In caso di connessioni multi-hop (con più nodi intermedi), il bitrate complessivo è dato dal bitrate del collegamento
più lento. Se si trasmette a bitrate superiori, si verificano fenomeni di congestione e perdita di pacchetti in quanto
il collegamento più lento non riesce a gestire il flusso di pacchetti in arrivo.

\subsection{Metriche di ritardo}
\subsubsection*{Componenti macroscopiche del ritardo della rete}
Le metriche di ritardo misurano il tempo impiegato dai dati per viaggiare da un nodo sorgente a un nodo destinazione
attraverso la rete. Di seguito i principali tipi di ritardo a livello macroscopico di rete:
\begin{itemize}
	\item \textbf{ritardo punto-punto} \(d_{p2p}\): tempo che intercorre tra l'inizio della trasmissione di un pacchetto
	in un nodo e la ricezione completa dello stesso pacchetto in un nodo adiacente, a volte è indicato anche come ritardo
	nodale \(d_\text{nodal}\)
	\item \textbf{ritardo end-to-end} \(d_{e2e}\): tempo che intercorre tra l'inizio della trasmissione di un pacchetto
	in un nodo e la ricezione completa dello stesso pacchetto nel nodo destinazione, è il risultato della somma di tutti
	i ritardi nodali lungo il percorso del pacchetto attraverso la rete
	\[d_{e2e} = \sum_{i = 1}^K d_\text{nodal,\,i} \qquad\qquad d_{e2e} = K \cdot d_\text{nodal} \text{ se i link sono identici}\]
	\item \textbf{round trip time} \(d_\text{RTT}\): tempo che intercorre tra l'inizio della trasmissione di un pacchetto in
	un nodo e la ricezione completa della risposta al pacchetto nello stesso nodo, è il risultato della somma dei ritardi
	end-to-end della trasmissione del pacchetto di richiesta e della risposta
	\[d_\text{RTT} = 2 \cdot d_{e2e}\]
\end{itemize}

\subsubsection*{Componenti del ritardo nodale}
Il ritardo nodale \(d_\text{nodal}\) si può scomporre nelle seguenti componenti:
\[d_\text{nodal} = d_\text{proc} + d_\text{queue} + d_\text{trans} + d_\text{prop}\]
\begin{itemize}
	\item \textbf{ritardo di processing} \(d_{proc}\): tempo necessario al nodo per elaborare il pacchetto, può essere
	dovuto a error checking, routing, ecc.; in genere è trascurabile siccome l'elaborazione è eseguita in hardware dedicati
	come ASIC, SmartNIC e DPU dedicate
	\item \textbf{ritardo di accodamento} \(d_{queue}\): tempo che il pacchetto trascorre nel buffer in attesa di essere
	trasmesso sul link, dipende dal traffico di rete e può essere molto variabile; può essere parzialmente ridotto con
	l'uso di sistemi di scheduling, priorità e controllo di ammissione dei pacchetti
	\item \textbf{ritardo di trasmissione} \(d_{trans}\): tempo che intercorre tra l'invio del primo bit e l'invio dell'ultimo
	bit del pacchetto sul link, equivale al rapporto tra la dimensione del pacchetto e il bitrate del link; può essere
	ridotto aumentando il bitrate del link; è detto anche tempo di forwarding o di inoltro
	\[d_\text{trans} = \frac{L}{R} = \frac{|PDU_\text{PHY}|}{S_\text{PHY}} \qquad \text{con } L \text{ dimensione del pacchetto e } R \text{ bitrate del link}\]
	\item \textbf{ritardo di propagazione} \(d_{prop}\): tempo necessario al pacchetto per propagarsi lungo il link tra due
	nodi adiacenti, equivale al rapporto tra la distanza tra i due nodi e la velocità di propagazione del segnale sul mezzo
	fisico; solitamente è trascurabile tranne in caso di link a lunga distanza come collegamenti satellitari o transoceanici;
	non può essere ridotto siccome non si può controllare la velocità di propagazione del segnale, al limite si può ridurre
	la distanza tra i due nodi
	\[d_\text{prop} = \frac{d}{v} \qquad \text{con } d \text{ distanza tra i due nodi e } v \text{ velocità di propagazione}\]
\end{itemize}

\subsubsection*{Jitter}
Il jitter è la varianza del ritardo di trasmissione. Risulta critico per applicazioni in cui si hanno vincoli di latenza
stringenti (es. real-time come videoconferenze, gaming online, telemedicina, ecc.) oppure per applicazioni in cui si richiede
un flusso continuo e costante di dati (es. streaming video o audio). Il jitter può essere mitigato con l'uso di un buffer
a discapito di una maggiore latenza complessiva.

\subsection{Prodotto banda-ritardo (BDP) e ritrasmissione dei pacchetti}
\subsubsection*{Prodotto banda-ritardo}
Il prodotto banda-ritardo, o Bandwidth-Delay Product (BDP), indica il numero di bit o pacchetti che possono essere
simultaneamente in transito attraverso una connessione di rete tra due nodi. Si calcola come il prodotto tra il throughput
a regime stazionario \(S\) per il ritardo RTT \(d_\text{RTT}\) (invio del pacchetto e acknowledgment della ricezione):
\[BDP_\text{bit} = S \cdot d_\text{RTT} \qquad\qquad BDP_\text{pkt} = \frac{BDP_\text{bit}}{|PDU_\text{PHY}|}\]

In caso di connessioni multi-hop, il BDP complessivo è calcolato in funzione del bitrate complessivo della connessione,
ovvero quello del collegamento più lento. Il ritardo RTT risulta essere la somma dei ritardi RTT di tutti i collegamenti
della connessione.

\subsubsection*{Ritrasmissione dei pacchetti}
Per assicurarsi che un pacchetto sia stato correttamente ricevuto dal nodo destinatario, il nodo sorgente deve attendere un
acknowledgment (ACK) di conferma. Se l'ACK non arriva entro un certo tempo di timeout, il pacchetto viene considerato
perso o danneggiato ed eventualmente dovrà essere ritrasmesso. Per questo motivo serve mantenere un buffer di trasmissione
per salvare i pacchetti inviati in attesa di ricevere l'ACK, per averli pronti per un'eventuale ritrasmissione.

La dimensione del buffer di trasmissione è pari al numero di pacchetti che possono essere simultaneamente in transito
e di cui si deve attendere l'ACK, ovvero esattamente il prodotto banda-ritardo.

In caso di connessioni multi-hop, l'ACK viene inviato solo dal nodo destinatario finale, quindi il buffer di trasmissione
deve garantire la memorizzazione di tutti i pacchetti in transito lungo l'intero percorso della connessione.

\subsubsection*{Prodotto banda-ritardo per Long Fat Networks (LFT) e connessioni multi-hop}
Il prodotto banda-ritardo è particolarmente importante per le connessioni su Long Fat Networks (LFN), ovvero reti con
elevata larghezza di banda e grande latenza come backbone transcontinentali, connessioni tra datacenter, connessioni
satellitari e edge cloud.

In questi casi la latenza risulta essere molto elevata, di conseguenza anche il BDP risulta molto elevato e si richiede
un buffer di trasmissione molto grande. Il problema è che le dimensioni del buffer sono limitate dal protocollo di
trasporto. Se la dimensione massima non è sufficiente (è minore del BDP), non è possibile inviare altri pacchetti
senza prima aver ricevuto l'ACK per quelli contenuti nel buffer, causando un calo del throughput della connessione.

\subsection{Metriche di affidabilità}
\subsubsection*{Packet error rate - PER}
Il packet error rate (\(PER\)) equivale al rapporto in percentuale tra il numero di pacchetti persi o danneggiati e il numero
totale di pacchetti trasmessi. Indica la probabilità che un pacchetto venga perso o danneggiato durante la trasmissione.

\subsubsection*{Dalla probabilità d'errore sul bit al packet error rate}
Un canale di trasmissione punto-punto per bit viene modellato con un canale binario simmetrico (BSC) con probabilità
d'errore \(\varepsilon\) sul bit. Si assume che gli errori sul bit siano indipendenti e identicamente distribuiti.
È possibile passare dalla probabilità d'errore sul bit \(\varepsilon\) al packet error rate \(PER\) sul pacchetto
di dimensione \(L\) bit con la seguente formula:
\[PER = 1 - (1 - \varepsilon)^L \qquad\quad \text{con } (1 - \varepsilon)^L = \begin{array}{c} \text{\small probabilità di ricevere tutti gli } L \text{\small ~bit} \\ \text{\small che formano il pacchetto senza errori} \end{array}\]
È utile definire la probabilità \(P(l)\) di ricevere \(l\) bit errati su \(L\) bit trasmessi:
\[P(l) = \binom{L}{l} \varepsilon^l \; (1 - \varepsilon)^{L-l} \qquad\quad \text{con } \begin{array}{c c} \displaystyle \binom{L}{l} = \text{\small \# modi diversi di scegliere i bit errati}\end{array}\]

\subsubsection*{Codifica di canale per rilevamento e correzione degli errori}
Per poter rilevare e correggere gli errori sui bit introdotti dal canale, si aggiungono dei bit di ridondanza secondo
le tecniche di codifica di canale. Esistono due tipi di codifica di canale:
\begin{itemize}
	\item \textbf{codici a blocco}: si partiziona il flusso di bit in blocchi di \(k\) bit a cui si aggiungono \(r\) bit
	di ridondanza ottenendo blocchi in uscita di \(n = k + r\) bit che vengono trasmessi sul canale; si definisce il tasso
	di codifica \(R = k/n < 1\) come fattore di riduzione del throughput dovuto all'overhead dei bit di ridondanza

	\item \textbf{codici convoluzionali}: lavorano con lo stesso principio, ma l'uscita ad un certo istante di tempo dipende
	anche dai dati trasmessi agli istanti di tempo precedenti (è dotato di memoria)
\end{itemize}

\noindent
Attualmente si utilizzano le seguenti codifiche di canale:
\begin{itemize}
	\item \textbf{LDPC} o Low-Density Parity-Check: codici a blocco sparso con decodifica iterativa, utilizzati per dati
	utente trasmessi ad alta velocità (5G, Wi-Fi 7/8, Ottica 400G/800G)
	\item \textbf{polar codes}: codici basati su canali sintetici e decodifica successive cancellation, utilizzati per
	canali di controllo e pacchetti brevi (5G, 6G)
	\item \textbf{Reed-Solomon}: codici algebrici a blocchi su campi finiti, utilizzati per correzione di errori burst
	su link ottici o storage
	\item \textbf{turbo codes}: codici convoluzionali concatenati in parallelo, utilizzati nei sistemi legacy e in alcuni
	uplink 4G/5G
\end{itemize}

\subsubsection*{Rilevamento di errori e controllo di parità}
I codici a controllo di parità consistono nell'aggiungere un bit di parità \(p\) al blocco di \(k\) bit da trasmettere,
ottenendo un blocco di \(n = k + 1\) bit. Il bit di parità è calcolato come la somma modulo 2 dei \(k\) bit del blocco,
ovvero \(p=0\) se se il numero di bit a 1 è pari, \(p=1\) se il numero di bit a 1 è dispari.

Questa struttura permette di rilevare fino ad 1 errore sul blocco di \(n\) bit, ma non permette di correggerlo. Si definisce
la capacità di rilevamento errori su un flusso continuo pari a  \(1/n\) che corrisponde anche all'overhead introdotto.
È possibile aumentare la capacità di rilevamento errori aggiungendo più bit di parità, ma aumenterà di conseguenza anche
l'overhead e quindi il calo del throughput.

\subsubsection*{Correzione di errori e codice di Hamming(n,k)}
I codici di Hamming sono codici a blocco che permettono di rilevare e correggere errori sui bit. Il codice di Hamming(\(n\),\(k\))
aggiunge \(r = n - k\) bit di ridondanza ogni \(k\) bit di dati, ottenendo blocchi di \(n\) bit. Riesce a rilevare fino a
2 errori sul blocco e a correggere fino a 1 errore sul blocco. Un esempio comune è il codice di Hamming(7,4) che aggiunge
3 bit di ridondanza ogni 4 bit di dati, con blocchi di 7 bit.

\subsubsection*{Errori in burst e interleaving}
Si parla di errori in burst quando si hanno più errori consecutivi su un blocco di dati. Sono tipici dei canali di
trasmissione reali in cui si hanno fenomeni di fading, interferenze o rumore impulsivo che si protraggono nel tempo
e di conseguenza coinvolgono più bit consecutivi.

Il problema degli errori in burst si verifica quando si utilizzano codici a blocco, in quanto tutti gli errori si concentrano
su uno o pochi blocchi, rendendone impossibile il rilevamento e la correzione. Per risolvere questo problema si utilizza la
tecnica di interleaving che consiste nel rimescolare i bit di più blocchi in modo da distribuire gli errori e diminuire il
numero di errori locale di ciascun blocco.

Ad esempio è possibile comporre una matrice di \(x\) righe composte dagli \(n\) bit di \(x\) blocchi consecutivi, e
trasmettere i bit leggendoli per colonna. Lo svantaggio è che si introduce un ritardo di trasmissione in quanto per
decodificare il primo blocco è necessario attendere la ricezione di tutti gli \(x\) blocchi.

\subsubsection*{Compromessi della codifica di canale}
La codifica di canale comporta alcuni svantaggi:
\begin{itemize}
	\item calo del bitrate pari al tasso di codifica \(R\) a causa dell'overhead introdotto dai bit di ridondanza
	\item aumento della complessità computazionale per l'encoding e il decoding dei blocchi di dati, maggiore è la capacità
	di rilevamento e correzione errori, maggiore è la complessità computazionale
	\item aumento del ritardo di trasmissione a causa dell'interleaving e delle codifiche convoluzionali
\end{itemize}

\subsubsection*{Congestione, packet loss rate e packet delivery ratio}
Prima di essere inoltrati su un canale, i pacchetti vengono memorizzati in un buffer di trasmissione. Se il buffer si riempie
a causa di un eccessivo traffico di rete, non è più possibile memorizzare i pacchetti in arrivo e si verifica un buffer overflow.
I pacchetti in arrivo devono venire scartati secondo una delle possibili politiche di gestione del buffer come le seguenti:
\begin{itemize}
	\item drop tail policy: se c'è un buffer overflow, si scartano tutti i pacchetti in arrivo
	\item early dropping casuale: quando l'occupazione del buffer supera una certa soglia (es. 2/3), si scartano casualmente
	alcuni pacchetti in arrivo
\end{itemize}
In caso di perdita del pacchetto, questo può essere ritrasmesso dal nodo sorgente, dall'ultimo nodo intermedio o non essere
ritrasmesso affatto, in base al protocollo di trasporto utilizzato.

Si definisce il packet loss rate (\(PLOSS\)) è definito come il rapporto in percentuale tra il numero di pacchetti persi
e il numero totale di pacchetti trasmessi. Indica la probabilità che un pacchetto venga perso durante la trasmissione.
Spesso viene utilizzato in modo intercambiabile con il packet error rate o packet drop rate (\(PER\)).

Come metrica complementare, di definisce anche il packet delivery ratio (\(PDR\)) come il rapporto in percentuale tra il
numero di pacchetti ricevuti correttamente e il numero totale di pacchetti trasmessi e si calcola come \(PDR = 1 - PLOSS\).

Si nota che le metriche di affidabilità \(PER\), \(PLOSS\) e \(PDR\) dipendono fortemente dal livello di rete in cui vengono
misurate e dai protocolli implementati, in quanto dipende da come vengono gestiti il rilevamento, la correzione d'errori e
la ritrasmissione dei pacchetti.
